\documentclass[aspectratio=169,11pt]{beamer}
% Shared Beamer style for practice contest solution walkthroughs.
\usepackage[utf8]{inputenc}
\usepackage[T1]{fontenc}
\usepackage{lmodern}
\usepackage{amsmath,amssymb}
\usepackage{xcolor}
\usepackage{hyperref}
\usepackage{listings}

\definecolor{CPNavy}{HTML}{172033}
\definecolor{CPBlue}{HTML}{2563EB}
\definecolor{CPGreen}{HTML}{16A34A}
\definecolor{CPOrange}{HTML}{F97316}
\definecolor{CPGray}{HTML}{64748B}
\definecolor{CPLight}{HTML}{F8FAFC}
\definecolor{CPCodeBg}{HTML}{F1F5F9}
\definecolor{CPCodeRule}{HTML}{CBD5E1}

\usetheme{default}
\usefonttheme{professionalfonts}
\setbeamertemplate{navigation symbols}{}
\setbeamersize{text margin left=0.65cm,text margin right=0.65cm}
\setbeamertemplate{itemize item}{\color{CPBlue}$\blacktriangleright$}
\setbeamertemplate{itemize subitem}{\color{CPGray}$\triangleright$}

\setbeamercolor{normal text}{fg=CPNavy,bg=white}
\setbeamercolor{frametitle}{fg=white,bg=CPNavy}
\setbeamercolor{title}{fg=white,bg=CPNavy}
\setbeamercolor{subtitle}{fg=CPLight,bg=CPNavy}
\hypersetup{colorlinks=true,urlcolor=CPBlue}

\setbeamertemplate{frametitle}{%
  \nointerlineskip%
  \begin{beamercolorbox}[wd=\paperwidth,ht=0.88cm,dp=0.22cm,leftskip=0.55cm,rightskip=0.35cm]{frametitle}
    \usebeamerfont{frametitle}\insertframetitle\hfill\small\insertframenumber/\inserttotalframenumber
  \end{beamercolorbox}%
}

\setbeamertemplate{footline}{%
  \leavevmode%
  \hbox{%
    \begin{beamercolorbox}[wd=.58\paperwidth,ht=2.4ex,dp=1ex,leftskip=.5cm]{author in head/foot}%
      \color{CPGray}\insertshorttitle
    \end{beamercolorbox}%
    \begin{beamercolorbox}[wd=.42\paperwidth,ht=2.4ex,dp=1ex,rightskip=.5cm plus1fil]{date in head/foot}%
      \hfill\color{CPGray}\insertshortauthor
    \end{beamercolorbox}%
  }%
}

\setbeamertemplate{title page}{%
  \vspace*{1.25cm}
  \begin{beamercolorbox}[wd=\paperwidth,sep=0.65cm,leftskip=0.15cm,rightskip=0.15cm]{title}
    {\color{white}\Huge\bfseries\inserttitle\par}
    \vspace{0.35cm}
    {\color{CPLight}\Large\insertsubtitle\par}
    \vspace{0.6cm}
    {\color{CPLight}\large\insertauthor\par}
  \end{beamercolorbox}
  \vfill
}

\newcommand{\sectiontag}[1]{\textbf{\color{CPBlue}#1}}
\newenvironment{tightitemize}{\begin{itemize}\small\setlength{\itemsep}{0.18cm}}{\end{itemize}}
\lstdefinestyle{cpcode}{
  language=Python,
  basicstyle=\ttfamily\tiny,
  keywordstyle=\color{CPBlue}\bfseries,
  commentstyle=\color{CPGray}\itshape,
  stringstyle=\color{CPGreen},
  backgroundcolor=\color{CPCodeBg},
  frame=single,
  framerule=0.25pt,
  rulecolor=\color{CPCodeRule},
  showstringspaces=false,
  columns=fullflexible,
  keepspaces=true,
  breaklines=true,
  breakatwhitespace=false,
  tabsize=4,
  xleftmargin=0.08cm,
  xrightmargin=0.08cm,
  aboveskip=0.08cm,
  belowskip=0.08cm
}


\title[unionfind]{Union-Find}
\subtitle{Day 1 solution walkthrough}
\author[Solution Deck]{Competitive Programming Bootcamp}
\date{}

\begin{document}

\begin{frame}[plain]
  \titlepage
\end{frame}

% BEGIN GENERATED STATEMENT INTERPRETATION
\begin{frame}{Interpret the Word Problem}
\small
\sectiontag{Statement excerpts:}
\begin{tightitemize}
  \item \textit{``the sets containing \$a\$ and \$b\$ are joined''}
  \item \textit{``asking whether \$a\$ and \$b\$ belong to the same set''}
\end{tightitemize}

\vspace{0.18cm}
\sectiontag{Programming interpretation:}
\begin{tightitemize}
  \item Maintain a partition of N elements under merge operations.
  \item Answer each query by checking whether two elements have the same representative.
  \item Path compression and rank keep a million operations feasible.
\end{tightitemize}
\end{frame}
% END GENERATED STATEMENT INTERPRETATION







\begin{frame}{Problem Model}
\small
\sectiontag{Textbook connection:} CP4 2.4.2 Union-Find Disjoint Sets

\vspace{0.25cm}
\sectiontag{Core technique:} Disjoint set union

\vspace{0.25cm}
\begin{tightitemize}
  \item The problem maintains dynamic connectivity among numbered elements.
  \item Operation '=' joins two sets.
  \item Operation '?' asks whether two elements currently share a set.
\end{tightitemize}
\end{frame}

\begin{frame}{How to Recognize the Approach}
\begin{tightitemize}
  \item The graph only gains connections; it never deletes them.
  \item Connectivity queries after unions are exactly the union-find use case.
  \item Path compression and union by rank keep operations effectively constant.
\end{tightitemize}
\end{frame}

\begin{frame}{Algorithm}
\begin{tightitemize}
  \item Initialize parent[i] = i, rank[i] = 0, and size[i] = 1.
  \item findSet walks parent pointers to the representative and compresses the path.
  \item isSameSet compares the two representatives.
  \item unionSet attaches the lower-rank root under the higher-rank root.
  \item Process commands and buffer yes/no answers.
\end{tightitemize}
\end{frame}

\begin{frame}{Walkthrough of the Key State}
\begin{tightitemize}
  \item Representatives name the current connected component.
  \item Path compression rewires visited nodes directly to the root.
  \item Union by rank prevents long chains from forming.
\end{tightitemize}
\end{frame}

\begin{frame}{Why the Algorithm Is Correct}
\begin{tightitemize}
  \item Initially every element is alone, matching the problem state.
  \item Each union merges exactly the two components containing its endpoints.
  \item Two elements are connected exactly when their representatives are equal.
\end{tightitemize}
\end{frame}

\begin{frame}{Complexity and Implementation Notes}
\small
\sectiontag{Complexity:} O((N + Q) alpha(N)) time and O(N) memory; alpha(N) is effectively constant.

\vspace{0.3cm}
\begin{tightitemize}
  \item The iterative find avoids Python recursion-limit problems.
  \item Buffer output because many query answers may be printed.
  \item The extra size and set-count fields are useful library features, even if not all are needed here.
\end{tightitemize}
\end{frame}



% BEGIN GENERATED CODE WALKTHROUGH
\begin{frame}[fragile]{Code: Build the DSU Arrays}
\scriptsize
\sectiontag{Source lines:} \texttt{unionfind.py:52--82}

\vspace{0.08cm}
\begin{columns}[T,totalwidth=\textwidth]
\begin{column}{0.58\textwidth}
\begin{lstlisting}[style=cpcode]
def __init__(self, numItems: int, numSets: int = None) -> None:
    if numItems < 1:
        raise ValueError("numItems needs to be greater or equal to 1")
    elif numSets is not None and numItems < numSets:
        raise ValueError("numSets must not be greater than numItems")

    self._numSet = numItems
    self._parents = list(range(numItems))
    self._ranks = [0] * numItems
    self._sizes = [1] * numItems
    if numSets is None:
        numSets = numItems
    else:
        while self._numSet > numSets:
            self.unionSet(random.randrange(numItems), random.randrange(numItems))
\end{lstlisting}
\end{column}
\begin{column}{0.38\textwidth}
\sectiontag{What this chunk does}
\begin{tightitemize}
  \item Each element begins as the parent of itself.
  \item rank supports shallow merges; size supports optional set-size queries.
  \item \_numSet tracks how many disjoint components remain.
\end{tightitemize}
\end{column}
\end{columns}
\end{frame}

\begin{frame}[fragile]{Code: Find with Compression}
\scriptsize
\sectiontag{Source lines:} \texttt{unionfind.py:112--143}

\vspace{0.08cm}
\begin{columns}[T,totalwidth=\textwidth]
\begin{column}{0.58\textwidth}
\begin{lstlisting}[style=cpcode]
def findSet(self, element: int) -> int:
    children = []
    toFind = element

    while toFind != self._parents[toFind]:
        children.append(toFind)
        toFind = self._parents[toFind]

    for c in children:
        self._parents[c] = toFind

    return toFind
\end{lstlisting}
\end{column}
\begin{column}{0.38\textwidth}
\sectiontag{What this chunk does}
\begin{tightitemize}
  \item The loop walks parent pointers until it reaches the root representative.
  \item Visited children are stored so their parent can be rewritten directly to the root.
  \item Future find operations on those nodes become much faster.
\end{tightitemize}
\end{column}
\end{columns}
\end{frame}

\begin{frame}[fragile]{Code: Check Same Set}
\scriptsize
\sectiontag{Source lines:} \texttt{unionfind.py:145--155}

\vspace{0.08cm}
\begin{columns}[T,totalwidth=\textwidth]
\begin{column}{0.58\textwidth}
\begin{lstlisting}[style=cpcode]
def isSameSet(self, element1: int, element2: int) -> bool:
    return self.findSet(element1) == self.findSet(element2)
\end{lstlisting}
\end{column}
\begin{column}{0.38\textwidth}
\sectiontag{What this chunk does}
\begin{tightitemize}
  \item Two elements are in the same set exactly when their representatives match.
  \item Calling findSet also compresses paths along the way.
\end{tightitemize}
\end{column}
\end{columns}
\end{frame}

\begin{frame}[fragile]{Code: Union by Rank}
\scriptsize
\sectiontag{Source lines:} \texttt{unionfind.py:157--182}

\vspace{0.08cm}
\begin{columns}[T,totalwidth=\textwidth]
\begin{column}{0.58\textwidth}
\begin{lstlisting}[style=cpcode]
def unionSet(self, element1: int, element2: int) -> None:
    s1 = self.findSet(element1)
    s2 = self.findSet(element2)

    if s1 == s2:
        return None

    if self._ranks[s1] > self._ranks[s2]:
        s1, s2 = s2, s1
    elif self._ranks[s1] == self._ranks[s2]:
        self._ranks[s2] += 1

    self._parents[s1] = s2

    self._numSet -= 1
    self._sizes[s2] += self._sizes[s1]
\end{lstlisting}
\end{column}
\begin{column}{0.38\textwidth}
\sectiontag{What this chunk does}
\begin{tightitemize}
  \item The method first finds both roots, because only roots should be linked.
  \item The smaller-rank root is attached below the larger-rank root.
  \item When ranks tie, one rank increases to reflect the possible height growth.
\end{tightitemize}
\end{column}
\end{columns}
\end{frame}

\begin{frame}[fragile]{Code: Process Commands}
\scriptsize
\sectiontag{Source lines:} \texttt{unionfind.py:210--230}

\vspace{0.08cm}
\begin{columns}[T,totalwidth=\textwidth]
\begin{column}{0.58\textwidth}
\begin{lstlisting}[style=cpcode]
lines = sys.stdin.read().strip().splitlines()
n = int(lines[0].split()[0])
output = []

ufds = UnionFind(n)

for command in lines[1:]:
    line = command.split()
    op = line[0]
    e1 = int(line[1])
    e2 = int(line[2])

    match op:
        case "?":
            output.append("yes\n" if ufds.isSameSet(e1, e2) else "no\n")
        case "=":
            ufds.unionSet(e1, e2)

sys.stdout.write("".join(output))
\end{lstlisting}
\end{column}
\begin{column}{0.38\textwidth}
\sectiontag{What this chunk does}
\begin{tightitemize}
  \item The input is read once, then commands are interpreted line by line.
  \item A question command appends yes or no based on isSameSet.
  \item An equals command merges the two sets with unionSet.
\end{tightitemize}
\end{column}
\end{columns}
\end{frame}
% END GENERATED CODE WALKTHROUGH

\begin{frame}{Source}
\small
\sectiontag{Solution file:} \texttt{practice\_contests/day\_1/unionfind.py}

\vspace{0.3cm}
\sectiontag{Problem:} \url{https://open.kattis.com/problems/unionfind}

\vspace{0.45cm}
Use this deck with the source code open beside it. The slides explain the reasoning; the code shows the exact input handling and output formatting.
\end{frame}

\end{document}
